% -----------------------------------------------------------------------------------
\section{What characterises ex-ante evaluation}\label{sec:exante}

The second half of the dichotomy claimed that menu coherence forces ex-ante evaluation. This section
shows that principal binary independence does force it, given symmetry, ordering, and monotonicity,
and that contraction does not.

\subsection{Binary independence: the Self-Esteem theorem}

\begin{theorem}[Relational Self-Esteem theorem]\label{thm:selfesteem}
On the domain of all finite real square matrices, a rule satisfies \ax{N}, \Bpr, and \ax{O} if and only if
there is a complete transitive relation $\succeq^*$ on $\Rl$ such that, for every $M$ and all $i,j$,
\[
  i\succeq_Mj\iff M_{ii}\succeq^*M_{jj}.
\]
\end{theorem}

\begin{proof}
Consider
\[
  H=\begin{pmatrix}p&r&p\\q&s&u\\p&t&p\end{pmatrix}.
\]
The principal block on $\{1,3\}$ is constant, so the transposition of $1$ and $3$ fixes it; by \ax{N} and
\Bpr{} the comparison of $1$ and $3$ is symmetric, and by completeness $1\sim3$. By transitivity, $1$ and
$3$ compare in the same way with $2$. By \Bpr{} the first comparison is that of the first index with the
second in $\bigl(\begin{smallmatrix}p&r\\q&s\end{smallmatrix}\bigr)$, and the second is that of the first
index with the second in the principal block on $(3,2)$, which is
$\bigl(\begin{smallmatrix}p&t\\u&s\end{smallmatrix}\bigr)$. Since $r,q,t,u$ are arbitrary, the comparison
of two options does not depend on the off-diagonal entries of their block, and by \Bpr{} and \ax{N} it is
a function of the ordered pair of diagonal values. Define $x\succeq^*y$ if the first index is weakly
preferred in $\operatorname{diag}(x,y)$. Completeness is inherited, and transitivity follows by applying
\Bpr{} to $\operatorname{diag}(x,y,z)$. Conversely, ranking by a weak order of the diagonal values satisfies
\ax{N}, \Bpr, and \ax{O}, and has nonempty maximal sets on finite menus.
\end{proof}

This is the relational form of the \enquote{Self-Esteem Theorem} of \textcite{podgorski2022}: if
the comparison of two options is a fixed function of the four values $V_X(X)$, $V_Y(Y)$, $V_X(Y)$,
$V_Y(X)$ and overall verdicts respect independence of irrelevant alternatives, the comparison
depends only on each option's value by its own lights, $V_X(X)$ and $V_Y(Y)$. Podgorski concludes
that among overall-value views only the evidential one respects that independence, and his
tournament decision theory escapes the theorem by producing no single overall value. Here the
common comparison function is supplied by \ax{N} and \Bpr, and the scalar overall value by \ax{O};
the relational version assumes no real-valued representation of $\succeq^*$.

\begin{corollary}\label{cor:three}
Adding \ax{F} leaves exactly three rules: increasing diagonal ranking, decreasing diagonal ranking,
and universal indifference. Adding \ax{Q} and \Mon{} leaves only increasing diagonal ranking.
\end{corollary}

\begin{proof}
Every pair $x>y$ is mapped to every other such pair by a positive affine transformation, so under \ax{F}
all unequal pairs are oriented alike; equal values tie by reflexivity and \ax{N}. \ax{Q} excludes
indifference. Under decreasing ranking, raising $M_{11}$ in the zero $2\times2$ matrix removes option $1$
from the choice set, violating the choice clause of \Mon.
\end{proof}

Since $M_{ii}=V_D(i)$,
\[
  \ax{N}+\ax{F}+\Bpr+\ax{O}+\ax{Q}+\Mon
\]
characterises ex-ante policy ranking under the stated matrix semantics. Under the structural axioms
\ax{N}, \ax{F}, \Bpr, \ax{O}, a single binary benchmark with unequal diagonal entries identifies which of the
three rules applies; no finite benchmark, however, certifies that the structural axioms themselves
hold everywhere.

\begin{proposition}[Exact domain threshold]\label{prop:threshold}
For domains consisting of all matrices up to a maximum size, the characterisation holds if and only if
matrices of size three are included.
\end{proposition}

\begin{proof}
With every matrix of sizes one to three, the proof of Theorem~\ref{thm:selfesteem} goes through, and
\Bpr{} propagates the conclusion to any larger sizes present. With sizes at most two, \Bpr{} is vacuous,
and ranking by $s_i=M_{ii}-M_{ij}$ satisfies \ax{N}, \ax{F}, \ax{O}, \ax{Q}, and \Mon; on
$\bigl(\begin{smallmatrix}1&100\\0&0\end{smallmatrix}\bigr)$ it chooses option $2$.
\end{proof}

Merely having some three-option problems is not enough; a sufficient richness condition is that the
domain contain all binary matrices, all the clone extensions $H$ above, and their relabellings (and,
for the affine conclusion, the relevant affine transforms). The argument also works on
$[0,1]$-valued matrices of sizes up to three, with \ax{F} interpreted internally.

\subsection{Contraction does not force the diagonal}

Can contraction replace binary independence? Under every menu semantics, the answer is no.

\begin{theorem}[A contraction-consistent family without diagonal optimality]\label{thm:lambda}
Write $d_i=M_{ii}$. For $\lambda\ge0$ let $i\to_\lambda j$ if $d_i-d_j>\lambda\lvert M_{ij}-M_{ji}\rvert$; choose
the options with no incoming edge, tied at the top, and rank the rejected options below them by their
diagonal entries. Every such rule satisfies \ax{N}, \ax{F}, \Cpr, \Ccol, \ax{O}, \ax{Q}, and \Mon. At
$\lambda=0$ it is increasing diagonal ranking; for every $\lambda>0$ it sometimes recommends an option
whose diagonal entry is not maximal.
\end{theorem}

\begin{proof}
Every edge strictly lowers the diagonal, so the edge relation is acyclic and the choice set is
nonempty. Whether $i\to_\lambda j$ depends only on $M_{ii},M_{jj},M_{ij},M_{ji}$, which survive both principal
and column restriction; so a chosen option acquires no incoming edge under restriction, which gives
both forms of contraction. Raising $d_i$ can only remove edges into $i$ and add edges out of $i$; hence
$i$'s chosen status can only improve, other options' chosen status can only worsen, and if $i$ stays
rejected its rank among the rejected improves, which gives \Mon. \ax{N} and \ax{F} are clear, since
$\lvert M_{ij}-M_{ji}\rvert$ scales with $a$ and is translation invariant. For $\lambda>0$, in
$\bigl(\begin{smallmatrix}1&0\\100/\lambda&0\end{smallmatrix}\bigr)$ there is no edge, and both options
are chosen.
\end{proof}

Even an ordered, neutral, affine-invariant, monotone, contraction-consistent rule, therefore, need not
maximise ex-ante value. The family also defeats finite certification: given finitely many benchmark
matrices, choose $\lambda>0$ so small that every option with non-maximal diagonal in every benchmark
receives an edge from a diagonal maximiser. The whole ranking then coincides with diagonal ranking on
every benchmark, while a new binary matrix with a large off-diagonal asymmetry relative to its
diagonal gap violates \ax{D}.

\begin{proposition}[Maximal cardinality]\label{prop:cardinality}
Let $\cont=\lvert\Rl\rvert$. There are $2^{\cont}$ rules satisfying \ax{N}, \ax{F}, \Cpr, \Ccol, \ax{O}, \ax{Q},
and \Mon.
\end{proposition}

\begin{proof}
For $E\subseteq\Rl$, put $i\to j$ if $d_i>d_j$, $s=M_{ji}-M_{ij}>0$, and either $d_i-d_j>s$, or $d_i-d_j=s$ and
$(d_j-M_{ij})/s\in E$; complete as in Theorem~\ref{thm:lambda}. Edges descend the diagonal and depend only
on principal blocks, so the choice set is nonempty and both contractions hold. Raising $d_i$ only
decreases the gaps $d_k-d_i$, so it never creates an edge into $i$ (at the boundary it may delete one),
and only increases the gaps $d_i-d_j$, so it never deletes an edge out of $i$; this gives \Mon. The
boundary parameter $(d_j-M_{ij})/s$ is invariant under positive affine maps. On
\[
  M_r=\begin{pmatrix}r+1&0\\1&r\end{pmatrix}
\]
we have $d_1-d_2=s=1$ and boundary parameter $r$, so the choice is $\{1\}$ if $r\in E$ and $\{1,2\}$
otherwise. Distinct sets $E$ give distinct rules, and $2^{\cont}$ is the number of all rules on finite real
matrices over a countable label universe.
\end{proof}

The class cannot, therefore, be parameterised by finitely many real numbers, or even countably many.

\subsection{Fixed-reference scores}

Suppose that only the score vector $v_a=\sum_jw_jM_{ja}$ is available, and that restriction deletes
coordinates. Even then rigidity does not return. The rule choosing the top $k\ge2$ scores, including
boundary ties, $G_k(v)=\{a:\#\{b:v_b>v_a\}<k\}$, with the chosen options tied above the others and the
others ranked by score, satisfies \ax{N}, \ax{F}, contraction, \ax{O}, \ax{Q}, and \Mon.

\begin{proposition}\label{prop:scores}
There are $2^{\cont}$ rules on finite score vectors satisfying \ax{N}, \ax{F}, contraction, \ax{O}, \ax{Q}, and
\Mon, already distinguished on vectors of length four.
\end{proposition}

\begin{proof}
Fix a set $E$ of normalised shapes of triples. For a candidate $j$ and a triple $T$ of other options,
let $g=(v_i-v_j)_{i\in T}$. Reject $j$ if some triple has $\sum g>0$, or has $\sum g=0$, $g\ne0$, and
sorted normalised shape in $E$; choose the rest, and complete the ranking as above. A score maximiser
is never rejected, since all its gaps are non-positive. Deletion removes witnesses, so contraction
holds. Raising $v_j$ lowers every gap sum for $j$, so a zero-sum rejection can only turn into a
negative sum, and raises the relevant sums for other candidates; this gives \Mon. Positive affine
maps preserve signs and normalised shapes. On $v=(0,-1,t,1-t)$ with $0<t<\tfrac12$, the option of
score $0$ has the single triple $(-1,t,1-t)$, of sum zero, and is rejected exactly when its shape lies
in $E$; these shapes are distinct for distinct $t$.
\end{proof}

With a fixed full-support $w$, raising $M_{jj}$ raises only $v_j$, so the matrix form of \Mon{} follows.
The upshot is that the dichotomy fails in both directions: contraction under any of the three
semantics is compatible with a rich family of non-ex-ante rules, and row-baseline-invariant
evaluation is compatible with contraction once the belief model is fixed. What does force ex-ante
evaluation is principal binary independence, which is a strong requirement: it imposes on every
comparison the beliefs that the two compared options would induce, and no others.
